\chapter{Human Bites}

\section*{Definition and Overview}
Human bites are wounds caused by human teeth, either from an intentional bite or, more commonly, from a punch to the mouth that cuts the knuckles ("fight bite"). Human mouths contain a wide range of bacteria, so bite wounds frequently become infected, and they require more careful treatment than most other wounds. The greatest risk is infection of the hand, which can spread to joints and tendons and cause permanent disability if not treated promptly. Cleaning, antibiotics, and review are the mainstays of management.

\section*{Epidemiology}
Human bites are common, particularly in children, where most bites come from other children and are relatively minor. In adults, bites often occur during fights or when caring for confused patients, and the hand is the most dangerous site. Up to a third of untreated human bite wounds become infected, and hand bites have the highest rate of serious complications. Infection of the hand from a fight bite can lead to joint and tendon infection requiring surgery and can cause permanent stiffness.

\section*{Aetiology and Pathophysiology}
\begin{itemize}
    \item \textbf{Mechanism:} teeth puncture the skin and introduce bacteria deep into tissues
    \item \textbf{Oral bacteria:} the human mouth contains more than 500 species, including streptococci, staphylococci, and Eikenella corrodens
    \item \textbf{Fight bite:} punching teeth cuts the knuckle; when the hand is opened, the wound seals over, trapping bacteria in the joint and tendon sheath
    \item \textbf{Pathophysiology:} bacteria multiply in the wound, causing cellulitis, abscess, and infection of joints (septic arthritis) and tendons (tenosynovitis)
    \item \textbf{Other risks:} transmission of hepatitis B, hepatitis C, and rarely HIV from infected blood
    \item \textbf{Tetanus:} possible if vaccination is not up to date
\end{itemize}

\section*{Clinical Features}
\begin{itemize}
    \item A bite mark, puncture wound, or laceration
    \item Pain, swelling, and redness developing over hours to days
    \item A hand wound near the knuckles, even if small
    \item Fever and feeling unwell with spreading infection
    \item Difficulty moving a finger or the hand
    \item Pus or discharge from the wound
    \item Deep wounds may damage tendons, nerves, or blood vessels
\end{itemize}

\section*{Differential Diagnosis}
\begin{itemize}
    \item Other animal bites (dog, cat) have different organisms but similar principles
    \item Uninfected bites, which need cleaning but not necessarily antibiotics
    \item The fight bite must be identified by history and examined with the hand in a fist
    \item Septic arthritis and tenosynovitis are complications to be excluded in hand bites
    \item The risk of infection determines whether antibiotics are given
\end{itemize}

\section*{When to Seek Medical Care}
All human bites that break the skin should be assessed by a doctor, and hand bites should be assessed urgently. Even small puncture wounds can become seriously infected.

\textbf{Call 000 (or local emergency services) for:}
\begin{itemize}
    \item Any bite to the hand, especially over a knuckle
    \item Spreading redness, severe pain, or swelling
    \item Fever or feeling unwell
    \item Difficulty moving a finger or the hand
    \item Deep wounds with heavy bleeding
    \item A bite with damage to tendons, nerves, or bone
\end{itemize}

\section*{Diagnosis and Investigations}
\begin{itemize}
    \item \textbf{History:} how the bite occurred, including whether it was a punch to the mouth
    \item \textbf{Examination:} wound depth, site, and signs of infection; the hand is examined in a fist position
    \item \textbf{Assessment of structures:} tendon, nerve, and joint function in hand bites
    \item \textbf{Wound swabs:} for culture when infection is established
    \item \textbf{X-ray:} to detect fracture or joint damage, especially in the hand
    \item \textbf{Blood tests:} for signs of systemic infection
\end{itemize}

\section*{Management}
\begin{itemize}
    \item \textbf{Cleaning:} immediate, thorough cleaning with soap and water or saline
    \item \textbf{Irrigation:} the wound is flushed out with large volumes of fluid
    \item \textbf{Debridement:} removal of damaged tissue by a doctor
    \item \textbf{Antibiotics:} a course of oral antibiotics is usually given for wounds that break the skin
    \item \textbf{Tetanus:} booster vaccination if not up to date
    \item \textbf{Wound closure:} most bites are left open to allow drainage; closure is reserved for selected wounds
    \item \textbf{Hand bites:} splinting, elevation, and urgent specialist review
    \item \textbf{HIV and hepatitis review:} post-exposure assessment where the biter is high risk
\end{itemize}

\section*{Complications}
\begin{itemize}
    \item Cellulitis and abscess formation
    \item Septic arthritis of the hand joints
    \item Tenosynovitis of the finger tendons
    \item Osteomyelitis (bone infection)
    \item Permanent stiffness and disability of the hand
    \item Scarring
    \item Transmission of blood-borne viruses, though the risk is low
\end{itemize}

\section*{Prognosis and Outcomes}
With prompt and thorough treatment, most human bites heal without complications. Bites elsewhere on the body usually do well with cleaning and antibiotics. Hand bites are the exception: they require urgent specialist treatment, and outcomes depend on how quickly the infection is controlled. With early surgery when needed, even serious hand infections usually recover, though some stiffness may persist.

\section*{Prevention and Self-Care}
\begin{itemize}
    \item Wash any bite wound immediately with soap and water
    \item Seek medical assessment for all bites that break the skin
    \item Treat hand bites as urgent
    \item Avoid punching and fighting, which cause the most dangerous bites
    \item Keep tetanus vaccination up to date
    \item Complete prescribed antibiotics fully
    \item Watch for signs of infection: spreading redness, pain, or fever
    \item Keep the wound clean and follow medical advice on wound care
\end{itemize}

\section*{Sources and Further Reading}
The following reputable sources provide further information for healthcare students and patients seeking reliable, current advice about human bites.

\begin{itemize}
    \item Healthdirect. \textit{Human bites: first aid and treatment.} https://www.healthdirect.gov.au/bites
    \item St John Ambulance Australia. \textit{Bite wound first aid.} https://stjohn.org.au/
    \item Therapeutic Guidelines. \textit{Bite wound management.} https://www.tg.org.au/
    \item NSW Health. \textit{Bites and stings fact sheets.} https://www.health.nsw.gov.au/
    \item Mayo Clinic. \textit{Human bites: first aid.} https://www.mayoclinic.org/
    \item MedlinePlus. \textit{Human bites.} https://medlineplus.gov/ency/article/000033.htm
\end{itemize}
